
Our work continues a long line of research on learning state abstractions for decision-making \cite{bertsekas1988adaptive,andre2002state,jong2005state,li2006towards,abel2017near,zhang2020learning}.
Most relevant are works that learn symbolic abstractions compatible with AI planners~\cite{lang2012exploration,jetchev2013learning,ugur2015bottom,asai2018classical,bonet2019learning,asai2020learning,ahmetoglu2020deepsym,umili2021learning}.
Our work is particularly influenced by \citet{pasula2007learning}, who use search through a concept language to invent symbolic state and action abstractions, and \citet{konidaris2018skills}, who discover symbolic abstractions by leveraging the initiation and termination sets of options that satisfy an abstract subgoal property.
The objectives used in these prior works are based on variations of auto-encoding, prediction error, or bisimulation, which stem from the perspective that the abstractions should \emph{replace} planning in the original transition space, rather than \emph{guide} it.

Recent works have also considered learning abstractions for multi-level planning, like those in the task and motion planning (TAMP)~\cite{gravot2005asymov,garrett2021integrated} and hierarchical planning~\cite{bercher2019survey} literature.
Some of these efforts consider learning symbolic action abstractions~\cite{zhuo2009learning,nguyen2017automated,silver2021learning,aineto2022comprehensive} or refinement strategies~\cite{chitnis2016guided,mandalika2019generalized,chitnis2019learning,wang2021learning,chitnis2021learning,ortiz2022structured}; our operator and sampler learning methods take inspiration from these prior works.
Recent efforts by \citet{loula2019discovering} and \citet{curtis2021discovering} consider learning both state and action abstractions for TAMP, like we do \cite{loula2019discovering,loula2020learning,curtis2021discovering}.
The main distinguishing feature of our work is that our abstraction learning framework explicitly optimizes an objective that considers downstream planning efficiency.

% More broadly, our work is situated within the larger research agenda of \emph{learning to plan}.
% Other efforts in this space include heuristic learning \cite{arfaee2011learning,silver2016mastering} (sometimes for TAMP \cite{shen2020learning,kim19learning,ltamppaxton}), learning to estimate the feasibility of abstract plans \cite{driess2020rss,noseworthy2021active}, and learning to generate reductions of planning models \cite{chitnis2020camps,silver2020planning}.
% Our efforts are also related to neuro-symbolic approaches that make use of symbolic planners, often in the context of hierarchical reinforcement learning \cite{grounds2005combining,yang2018peorl,kokel2021reprel}; we differ in that the agent knows the environment transition model in our setting, but is trying to learn abstractions that making planning under that model \emph{efficient}.
% A separate line of work studies learning symbolic planning models from plan traces~\cite{zhuo2010learning,yang2007learning,zhuo2013refining}; we are also interested in learning planning models, but in our setting, the data and environment are continuous, and our symbolic models are used as abstractions for efficient planning.
% Finally, inverse planning~\cite{baker2009action,ramirez2010probabilistic,zhi2020online} and inverse reinforcement learning~\cite{ng2000algorithms,paxton2016want,kurutach2018learning,jang2022bc} study the problem of leveraging demonstrations to infer cost functions and, in turn, policies. In our setting, we also learn from demonstrations, but we differ in that we assume goals are known; the unknowns are effective state and action abstractions for planning in held-out tasks.